% ============================================================
% ch14.tex —— 第 14 章 逼近的技艺：刘徽的刀
% ============================================================
\chapter{逼近的技艺：刘徽的刀}

公元 263 年，魏晋乱世，数学家刘徽给《九章算术》作注，在"圆田术"一节写下了一段将照亮后世一千七百年的话：

\begin{quote}\centering\itshape
割之弥细，所失弥少，割之又割，以至于不可割，则与圆周合体而无所失矣。
\end{quote}

——把圆切得越细，损失就越少；切了又切，切到不能再切，多边形就与圆周融为一体，毫无损失。这一章我们把这段话锻造成一件精密工具：\textbf{极限}。它是后面所有章节的发动机。

\section{刘徽怎么算圆周率}

\begin{kaiwei}
不许用"$\pi \approx 3.14$"这个现成答案。只许用尺规和算术，怎么\textbf{算}出圆周率？
\end{kaiwei}

刘徽的刀法分三步。\textbf{第一步}：在半径为 $1$ 的圆里画内接正六边形——它有个漂亮的性质：边长恰好等于半径 $1$（每条边对应的圆心角 $60^\circ$，与两条半径恰好拼成等边三角形）。于是正六边形周长 $= 6$，周长的一半 $= 3$。这就是"径一周三"（直径 $1$ 周长 $3$）的出处——中国古代长期拿 $3$ 当圆周率用，刘徽对此评价：太粗了。

\textbf{第二步：翻倍切一刀}。把六边形切成正十二边形。新边长怎么算？刘徽给出了递推公式（用勾股定理两次就能推，此处直接使用）：正 $n$ 边形边长 $l_n$ 翻倍后
\[
l_{2n} = \sqrt{\left(\tfrac{l_n}{2}\right)^2 + \left(1 - \sqrt{1 - \left(\tfrac{l_n}{2}\right)^2}\right)^2\,} .
\]
从 $l_6 = 1$ 出发，逐刀切下去：

\begin{center}
\begin{tabular}{ccc}
\toprule
正多边形 & 周长的一半（$\pi$ 的近似） & 与 $\pi$ 的差距 \\
\midrule
$6$ 边形 & $3.000000$ & 差 $0.1416$ \\
$12$ 边形 & $3.105829$ & 差 $0.0358$ \\
$24$ 边形 & $3.132627$ & 差 $0.0090$ \\
$48$ 边形 & $3.139350$ & 差 $0.0022$ \\
$96$ 边形 & $3.141032$ & 差 $0.0006$ \\
$\cdots$ & $\cdots\ \to\ 3.141593\cdots$ & $\to\ 0$ \\
\bottomrule
\end{tabular}
\end{center}

\begin{tikzpicture}[scale=1.15]
  \draw[thick] (0,0) circle (1.5);
  \draw[shenlan] (90:1.5) \foreach \a in {150,210,270,330,30} { -- (\a:1.5) } -- cycle;
  \draw[molv] (90:1.5) \foreach \a in {120,150,180,210,240,270,300,330,0,30,60} { -- (\a:1.5) } -- cycle;
  \fill (0,0) circle (0.8pt);
  \node[shenlan] at (2.2,0.95) {正六边形：周长率 $3.0$};
  \node[molv] at (2.35,-0.95) {正十二边形：$3.106$};
  \node at (0,-2.15) {割之弥细，所失弥少};
\end{tikzpicture}

\textbf{第三步：见证规律}。每切一刀，数字都\textbf{更大一点}（多边形把圆包得更满——每个新顶点都从弦的"凹处"顶到圆弧上），差距\textbf{大约缩小到四分之一}。刘徽自己写到 $96$ 边形，得到
\[
3.141024 < \pi < 3.142704
\]
（他连误差的上下界都夹出来了——"夹"这个动作马上会成为本章主角）。两百年后，祖冲之父子把这场割圆比赛推到正 $24576$ 边形，在没有任何计算器的年代得到 $3.1415926 < \pi < 3.1415927$，这个世界纪录保持了九百多年。

\section{核心对象：永不抵达、却越靠越近的一串数}

刘徽的表格里藏着本书最重要的一个抽象。抽出来看：
\[
3.0,\quad 3.1058,\quad 3.1326,\quad 3.1394,\quad 3.1410,\ \dots
\]
这是一串数（\textbf{数列}）。它有两个特征：
\begin{enumerate}
\item \textbf{单调上升}——每个数都比前一个大（多边形越切越包满）；
\item \textbf{有上界}——全都小于某个定数（比如 $3.1427$，圆周长本身当然是天花板）。
\end{enumerate}
单调上升且有上界的数列，必然越挤越紧地靠向某个数——它\textbf{永不抵达}（$96$ 边形还是多边形，不是圆），却\textbf{越靠越近}。那个"被靠"的数，就是这一串数的\Keyword{落脚点}，$\pi$。

第 13 章的"路标与目的地"在此完全兑现：$3.0, 3.1058, 3.1326, \dots$ 是路标串，$\pi$ 是目的地。\textbf{目的地不属于路标串}（没有哪个正多边形的周长率恰好等于 $\pi$），但路标串的"归宿"由它唯一决定。上一章 $0.999\cdots = 1$ 说的完全是同一件事：路标串 $0.9, 0.99, 0.999, \dots$ 的落脚点是 $1$。

\begin{faming}
\textbf{土名}：落脚点 $\to$ \textbf{正式名}：\textbf{极限}（limit）。数列 $a_1, a_2, a_3, \dots$ 收敛到 $A$，记作
\[
\lim_{n\to\infty} a_n = A ,
\]
读作"当 $n$ 跑向无穷时，$a_n$ 的落脚点是 $A$"。"$\lim$"是拉丁文 limes（界限）的缩写。
\end{faming}

\section{落脚点的严格规矩：一场必胜的比赛}

"越靠越近"是形容词，数学要把它铸成铁律。规矩设计成一场\textbf{两人对抗赛}：

\begin{dingyi}[收敛的比赛规则（$\varepsilon$--$N$ 语言的直白版）]
裁判宣布：数列 $a_n$ 收敛到 $A$，当且仅当——
\begin{quote}
\textbf{无论对手挑一个多小的误差 $\varepsilon$（零点一、千分之一、$10^{-100}$，随便挑），你总能报出一个项数 $N$，使得第 $N$ 项之后的全体成员都满足}
\[
|a_n - A| < \varepsilon .
\]
\textbf{你能永远应战，就判你赢（收敛成立）；只要有一个 $\varepsilon$ 你应不了战，就判输。}
\end{quote}
\end{dingyi}

用这场"要多近有多近"的比赛重审刘徽的表：

\textbf{第一回合}：对手挑 $\varepsilon = 0.01$。你翻表：$24$ 边形周长率 $3.1326$，与 $\pi$ 差 $0.0090 < 0.01$，而且之后的每一项差距更小。你报 $N = 24$ 边形那一行。应战成功。

\textbf{第二回合}：对手挑 $\varepsilon = 0.001$。表里 $96$ 边形差 $0.0006 < 0.001$，报 $N = 96$。成功。

\textbf{第三回合}：对手挑 $\varepsilon = 10^{-10}$。表里没有——但你手里有递推公式：每切一刀差距缩到约四分之一，从 $0.0006$ 出发再切 $17$ 刀（$4^{17} \approx 1.7\times10^{10}$），差距压到 $10^{-10}$ 以下。报 $N = 96\times2^{17}$ 边形。成功。

关键领悟：\textbf{对手无论挑多小，你都有应手}——不是因为你比对手快，是因为"翻倍"这台机器有无穷的后劲。\textbf{极限从来不是一个地点，是一场必胜的比赛}。谁能永远应战，谁的落脚点就存在。这套规则在大学课本里写作"$\forall\varepsilon>0, \exists N, \forall n>N: |a_n - A|<\varepsilon$"——符号不同，骨架一字不差。

现在还第 13 章的债。路标串 $a_n = 0.\underbrace{99\cdots9}_{n\text{个}}$ 的落脚点是 $1$ 吗？比赛验证：对手挑 $\varepsilon$（不妨设 $10^{-k}$），你报 $N = k$——第 $k$ 项之后的所有项与 $1$ 的差都不超过 $10^{-k} < \varepsilon$（差恰好是 $0.00\cdots01$，$k$ 位）。应战成功，且对一切 $\varepsilon$ 成立：$\lim a_n = 1$。\textbf{无穷位小数"$0.999\cdots$"的正式定义就是这场比赛的胜利}，论证一预支的信用（"无穷小数可以逐位运算"）到账：收敛数列的四则运算合法（和差的极限 $=$ 极限的和差），所以"错位相减"每一步都清白。

\section{实战：用比赛规则算三个极限}

光有规则不干活不行。极限计算的头号技巧一句话：\textbf{把"无穷大"的表达式，除成"零与常数"的组合}——无穷的问题瞬间变成算术。

\paragraph{例一}：$a_n = \dfrac{2n^2 + n}{n^2 + 3}$。先看几项：$0.75, 1.0, 1.1, 1.14, \dots$——趋势不明（爬得很慢）。上下同除以 $n^2$（上下同除一个数，分数不变，初中规则）：
\[
a_n = \frac{2 + \frac{1}{n}}{1 + \frac{3}{n^2}} .
\]
$n$ 极大时，$\tfrac1n$ 与 $\tfrac{3}{n^2}$ 是可压缩项：$n = 10^6$ 时它们分别是 $0.000001$ 和 $0.000000000003$——比赛里压成零。所以
\[
\lim_{n\to\infty} a_n = \frac{2+0}{1+0} = 2 .
\]
\textbf{慢，但方向从第一项就定了}——这就是极限眼看穿无穷远处的能力。

\paragraph{例二}：$b_n = \dfrac{5n}{n+1}$。同除 $n$：$\dfrac{5}{1 + \tfrac1n} \to \dfrac51 = 5$。

\paragraph{例三（技巧升级）}：$c_n = \sqrt{n^2+n} - n$。直接"代入无穷大"得 $\infty - \infty$——\textbf{禁区}（两个无穷大相减，胜负未知）。技巧：分子有理化（乘共轭）：
\[
c_n = \frac{(\sqrt{n^2+n} - n)(\sqrt{n^2+n} + n)}{\sqrt{n^2+n} + n} = \frac{n}{\sqrt{n^2+n} + n},
\]
再上下同除 $n$（注意 $\sqrt{n^2+n} \div n = \sqrt{1 + \tfrac1n}$）：
\[
c_n = \frac{1}{\sqrt{1+\frac1n} + 1} \;\longrightarrow\; \frac{1}{\sqrt1 + 1} = \frac12 .
\]
无穷大减无穷大，答案竟是老实的 $\tfrac12$——\textbf{禁区不是禁区，是"需要绕行的路段"}。有理化 $+$ 同除，是本书极限工具箱的双扳手。

\section{两块警示牌}

\textbf{警示牌一：不是所有数列都有落脚点。}$1, -1, 1, -1, \dots$ 永远摇摆：对手挑 $\varepsilon = 0.1$，你报任何 $N$，第 $N$ 项之后仍然既有 $1$ 又有 $-1$，它们与任何候选终点 $A$ 的距离不可能同时小于 $0.1$（$1$ 与 $-1$ 相距 $2$，两头够不着）。你永远应不了战——\textbf{发散}（摆动型）。数列 $1, 2, 3, \dots$ 也发散（膨胀型）。"\textbf{先问收敛，再算极限}"从此成为流程。

\textbf{警示牌二：数列收敛不等于"很快到达"。}$a_n = \tfrac{1000000}{n}$ 的落脚点是 $0$，但它前 $50$ 万项都比 $1$ 大——前半程看起来毫无收敛迹象。比赛的裁判只看"\textbf{最终}"（第 $N$ 项之后），不看开局。这也是为什么"$10^6$ 项的和是 $14.4$"这种数据（第 21 章的调和级数）会骗过直觉：发散的级数，前段完全可以伪装成收敛。

\begin{lishi}
刘徽，生卒年不详（约 225--295），山东邹平人，布衣数学家。他在注《九章算术》时创立"割圆术"，还独立提出十进小数、负数运算等概念，且明确写下了"析之弥细"的方法论自觉——他知道自己在发明方法，而不只是算题。祖冲之（429--500）与儿子祖暅在割圆术的基础上算到 $24576$ 边形，并给出密率 $\tfrac{355}{113} = 3.1415929\cdots$（分母一万以内最佳近似，好记：把 $113355$ 劈成两半）。祖氏父子的著作《缀术》唐代被列为算学教科书，北宋时失传——我们今天知道他们的成就，靠的是《隋书·律历志》短短几十字的记载。
\end{lishi}

\begin{dongshou}
\begin{itemize}
  \yisi 求落脚点：$\dfrac{3n}{n+1}$；$\dfrac{5n^2+2}{n^2+n}$；$0.7, 0.77, 0.777, \dots$（第三个先猜，再用比赛规则验证；它的极限写成无穷小数是 $0.777\cdots = \tfrac79$，除法竖式验证）。
  \yier 实战演练比赛规则：数列 $a_n = \tfrac{n}{n+1}$，落脚点 $1$。对手挑 $\varepsilon = 0.01$，你报哪个 $N$？（要 $|a_n - 1| = \tfrac{1}{n+1} < 0.01$，即 $n+1 > 100$，报 $N = 100$。）对手再挑 $10^{-6}$、$10^{-100}$，各报几？
  \yier 证明数列 $a_n = (-1)^n\cdot\dfrac{n}{n+1}$ 没有落脚点。（模仿警示牌一：它的偶数项奔 $1$、奇数项奔 $-1$，任何一个候选终点都会被 $\varepsilon = 1$ 击败。）
  \yier 用"同除 + 有理化"双扳手算：$\lim \dfrac{n^2+1}{3n^2+2n}$；$\lim\left(\sqrt{n+1} - \sqrt n\right)$。（第二个乘共轭 $\sqrt{n+1}+\sqrt n$，答案 $0$——但它\textbf{从下方}逼近 $0$，恒正。）
\end{itemize}
\end{dongshou}

\begin{shaonao}
用比赛规则证明"\textbf{落脚点唯一}"：若 $a_n$ 同时收敛到 $A$ 和 $B$（$A \ne B$），令 $d = |A - B| > 0$。取 $\varepsilon = \tfrac d2$，按规则存在 $N_1$ 使 $n > N_1$ 时 $a_n$ 离 $A$ 不超过 $\tfrac d2$，也存在 $N_2$ 使 $n > N_2$ 时离 $B$ 不超过 $\tfrac d2$。取 $n$ 同时超过两者，用"三角不等式"（两点经过中转不短于直线：$|A - B| \le |A - a_n| + |a_n - B|$）推出 $d \le \tfrac d2 + \tfrac d2 = d$……恰好取等？不——取 $\varepsilon = \tfrac d3$ 再走一遍，得 $d \le \tfrac{2d}3 < d$，矛盾。所以落脚点至多一个。（这道题的证明套路与第 7 章"单位元唯一"异曲同工：都是让两个候选者互相"演对方"，逼出矛盾。公理推理与极限推理，骨头是同一根。）
\end{shaonao}

\begin{shaonao}
刘徽的"夹"字诀升级：若三个数列处处满足 $b_n \le a_n \le c_n$，且外侧两个收敛到\textbf{同一个}极限 $L$，则中间的 $a_n$ 也收敛到 $L$（\textbf{夹逼定理}）。用它做一道刘徽式的题：求 $\lim \dfrac{\sin(1/n)}{1/n}$……不行，三角函数超纲。换一道纯代数的：求
\[
\lim_{n\to\infty} \left(\frac{1}{\sqrt{n^2+1}} + \frac{1}{\sqrt{n^2+2}} + \cdots + \frac{1}{\sqrt{n^2+n}}\right) .
\]
（提示：每一项都夹在 $\tfrac{1}{\sqrt{n^2+n}}$ 与 $\tfrac{1}{\sqrt{n^2+1}}$ 之间，整串 $n$ 项之和便夹在 $\dfrac{n}{\sqrt{n^2+n}}$ 与 $\dfrac{n}{\sqrt{n^2+1}}$ 之间；两个外侧数列上下同除 $n$，都奔 $1$。答案 $1$。夹逼定理是第 17 章"上下阶梯夹住面积"的正式许可证。）
\end{shaonao}

\begin{chengshi}
本章的比赛规则是"$\varepsilon$--$N$ 语言"的直白版，两者一一对应。"单调有界必收敛"（本节的直觉依据）在实数理论里是一条公理（确界原理）的推论，大学分析第一课会讲；收敛数列四则运算的严格证明（约一页）我们略去，使用规则本身。割圆递推公式的推导需要勾股定理两连击（可自行尝试：从圆心作边的心垂线），表格数值直接引用刘徽《九章算术注》原算。
\end{chengshi}

刀已磨好。下一章，用它剖开一个四百年前把伽利略逼到墙角的问题：\textbf{"瞬间"的速度，到底存不存在？}——你将亲手把"导数"重新发明一遍。
